\setcounter{section}{7}

  \zwischenueberschrift{Soal latihan}

\inputaufgabe
{}
{

Tunjukkan bahwa untuk dua titik
\mathl{A,B \neq N,S}{} pada bola satuan $S^2$, selisih jarak yang diambil dalam dua peta standar stereografis, yaitu
\mathkor {} {d( \alpha_1(A),\alpha_1(B))} {dan} {d( \alpha_2(A),\alpha_2(B))} {,}
dapat dibuat sekecil dan sebesar apa pun.

}
{} {}
Soal di atas menunjukkan bahwa secara umum kita tidak dapat memperoleh pengertian jarak yang bermakna pada suatu manifold melalui peta. Metrik alami pada bola satuan diperoleh dari metrik terinduksi pada $\R^3$ atau dari jarak geodesik; lihat
Aufgabe 7.16.

\inputaufgabe
{}
{

Tunjukkan bahwa suatu
\definitionsverweis {interval setengah terbuka}{}{}
\mathl{[a,b[}{} bukan
\definitionsverweis {manifold topologis}{}{}.

}
{} {}

\inputaufgabe
{}
{

Misalkan
\mathl{I=[0,1[}{} adalah
\zusatzklammer {setengah terbuka ke atas} {} {}
interval satuan dan $S^1$ adalah
\definitionsverweis {lingkaran satuan}{}{.}
Tunjukkan bahwa terdapat suatu
\definitionsverweis {pemetaan bijektif}{}{}
\definitionsverweis {kontinu}{}{}
\maabbdisp {f} { I } { S^1
} {}
tetapi
\mathkor {} {I} {dan} {S^1} {}
tidak
\definitionsverweis {homeomorfik}{}{}
.

}
{} {}

\inputaufgabegibtloesung
{}
{

Tunjukkan bahwa ketiga manifold berdimensi satu
\mathlistdisp {S^1 = \left\{ (x,y) \in \R^2  \mid   \betrag { (x,y) } = 1         \right\}} {} {\R} {dan} {\R \setminus \{0\}} {}
tidak homeomorfik satu sama lain secara berpasangan.

}
{} {}

\inputaufgabe
{}
{

Misalkan $M$ adalah
\definitionsverweis {grafik}{}{}
dari fungsi

\mavergleichskettedisp
{\vergleichskette
{ f(x)
}
{ =} {\begin{cases} \sin \frac{1}{x} \text{ untuk } x \neq 0\, , \\ 0 \text{ selain itu }\, , \end{cases}
}
{ } {
}
{ } {
}
{ } {
}
}
{}{}{}
yang diberikan oleh pemetaan
\maabb {f} {\R} {\R
} {.}
Tunjukkan bahwa $M$ bukan
\definitionsverweis {manifold topologis}{}{}.

}
{} {}

\inputaufgabe
{}
{

Tunjukkan bahwa sebuah
\definitionsverweis {bola terbuka}{}{}

\mavergleichskette
{\vergleichskette
{ U \left(   P ,r  \right)
}
{ \subseteq }{ \R^m
}
{  }{
}
{    }{
}
{   }{
}
}
{}{}{}
adalah $C^\infty$-\definitionsverweis {difeomorfik}{}{}
dengan $\R^m$.

}
{} {}

\inputaufgabegibtloesung
{}
{

Misalkan

\mathdisp {S= \left\{ P \in \R^3  \mid   \Vert { P } \Vert = 1         \right\}} {  }

adalah bola satuan. Untuk
\mathl{v=(a,b,c) \neq 0}{} berlaku

\mathdisp {E_v = \left\{ (x,y,z) \in \R^3  \mid  ax+by+cz = 0         \right\}} {  }

sebuah bidang melalui titik nol, yang mendefinisikan sebuah lingkaran besar
\zusatzklammer {sebuah \anfuehrung{khatulistiwa}{}} {} {}
dan dua hemisfer terbuka pada $S$.

a) Jelaskan, untuk
\mathl{v=(a,b,c) \neq 0}{}, lingkaran besar terkait dan kedua hemisfer dengan persamaan atau pertidaksamaan.

b) Tunjukkan bahwa $S$ tidak dapat ditutupi oleh tiga hemisfer terbuka.

}
{} {}

\inputaufgabe
{}
{

Tunjukkan bahwa permukaan silinder terbuka
\mathl{S^1 \times {]0,1[}}{} homeomorfik dengan
\mathl{S^1 \times \R}{,} dengan bidang yang ditusuk
\mathl{\R^2 \setminus \{(0,0)\}}{} dan dengan
\mathl{S^2 \setminus \{N,S\}}{}
\definitionsverweis {homeomorfik}{}{.}

}
{} {}

\inputaufgabe
{}
{

Misalkan
\mathl{]a,b[}{} sebuah
\definitionsverweis {interval terbuka}{}{}
dan
\maabbdisp {f} {]a,b[} {\R_{\geq 0}
} {}
sebuah
\definitionsverweis {fungsi kontinu}{}{.}
Misalkan $M$ adalah permukaan luar benda hasil rotasi terkait. Tunjukkan bahwa ketika

\mavergleichskette
{\vergleichskette
{ f
}
{ > }{ 0
}
{  }{
}
{    }{
}
{   }{
}
}
{}{}{}
ia merupakan manifold yang homeomorfik dengan selubung silinder terbuka. Tunjukkan lebih lanjut bahwa tidak ada manifold jika $f$ memiliki baik titik nol maupun nilai positif.

}
{} {}

\inputaufgabe
{}
{

Tunjukkan bahwa sebuah
\definitionsverweis {manifold topologis}{}{}
\definitionsverweis {terhubung}{}{}
tepat jika dan hanya jika
\definitionsverweis {terhubung lintasan}{}{}
.

}
{} {}

\inputaufgabe
{}
{

Misalkan $M$ sebuah
\definitionsverweis {manifold topologis}{}{}
dan misalkan

\mathl{M= \bigcup_{i \in I} U_i}{} sebuah
\definitionsverweis {tutupan terbuka}{}{}
dengan
\definitionsverweis {peta}{}{}
\maabbdisp {\alpha_i} {U_i } { V_i
} {}
dengan
\mathl{V_i \subseteq \R^n}{.} Untuk
\mathl{i,j \in I}{} definisikan
\maabbdisp {\varphi_{ij} = \alpha_j \circ (\alpha_i)^{-1}} {V_i \cap \alpha_i(U_i \cap U_j)  } { V_j \cap \alpha_j(U_i \cap U_j)
} {}
yang disebut
\definitionsverweis {pemetaan transisi}{}{.}
Tunjukkan bahwa untuk
\mathl{i,j,k \in I}{} apa yang disebut \stichwort {syarat koklus} {}
\zusatzklammer {pada himpunan bagian mana} {?} {}

\mavergleichskettedisp
{\vergleichskette
{ \varphi_{ij}
}
{ =} {\varphi_{kj} \circ \varphi_{ik}
}
{ } {
}
{ } {
}
{ } {
}
}
{}{}{}
berlaku.

}
{} {}

\inputaufgabe
{}
{

Periksa syarat koklus
\zusatzklammer {lihat
Aufgabe 7.11} {} {}
untuk bola satuan $M=S^2$ dan tiga proyeksi stereografis dari Kutub Utara, Kutub Selatan, dan
\mathl{P=(1,0,0)}{.}

}
{} {}

\inputaufgabegibtloesung
{}
{

Buat sebuah animasi yang menampilkan objek-objek geometri dari
Aufgabe 7.17.

}
{} {}

\inputaufgabe
{}
{

Perhatikan kertas kado. Dengan cara apa kertas tersebut dapat dipotong dan/atau direkatkan sehingga menghasilkan manifold berdimensi dua? Apakah tepi kertas harus menjadi bagian darinya atau tidak? Manifold mana yang dihasilkan yang terhubung, mana yang kompak? Bagaimana pita Möbius dibuat? Kemungkinan apa yang ada jika kita mengizinkan sejumlah hingga titik pengecualian, tempat tidak ada struktur manifold?

Terapkan teori ini untuk merancang kemasan yang seorisinal mungkin. Ikatlah kemasan tersebut dengan manifold berdimensi satu yang sesuai.

}
{} {}

  \zwischenueberschrift{Soal untuk dikumpulkan}

\inputaufgabe
{3}
{

Tentukan
\definitionsverweis {citra}{}{} dari
\definitionsverweis {lingkaran-lingkaran besar}{}{}
melalui kedua kutub pada
\definitionsverweis {bola satuan}{}{} di bawah
\definitionsverweis {proyeksi stereografis}{}{}
dari Kutub Utara.

}
{} {}

\inputaufgabe
{5}
{

Tunjukkan bahwa pada
\definitionsverweis {bola satuan}{}{}

\mavergleichskette
{\vergleichskette
{ K
}
{ \subset }{ \R^3
}
{  }{
}
{    }{
}
{   }{
}
}
{}{}{}
melalui korespondensi berikut ditentukan sebuah
\definitionsverweis {metrik}{}{} . Untuk

\mavergleichskette
{\vergleichskette
{ P,Q
}
{ \in }{ K
}
{  }{
}
{    }{
}
{   }{
}
}
{}{}{}
\mathl{d(P,Q)}{} adalah panjang lintasan penghubung
\zusatzklammer {yang lebih pendek} {} {} dari $P$ ke $Q$ pada
\definitionsverweis {lingkaran besar}{}{}
\zusatzklammer {perhatikan juga kasus

\mavergleichskette
{\vergleichskette
{ P
}
{ = }{ Q
}
{  }{
}
{    }{
}
{   }{
}
}
{}{}{}
dan $P,Q$
\definitionsverweis {antipodal}{}{}} {} {.}

}
{} {}

\inputaufgabe
{8}
{

Kita tetapkan dua titik $N=(0,0,1)$ dan $P=(1,0,0)$ pada
\definitionsverweis {bola satuan}{}{} $K$. Misalkan $G$ adalah garis penghubung dan $H$ adalah bidang tegak lurus terhadap $G$ melalui $N$. Pada $H$, buat lingkaran satuan berparameter $E$ dengan $N$ sebagai pusat. Tentukan, untuk $S \in E$, panjang lintasan
\zusatzklammer {yang lebih pendek} {} {} dari $N$ ke $P$ pada lingkaran yang ditentukan oleh irisan $K$ dengan bidang yang melalui
\mathkor {} {N,P} {dan} {S} {},
tersebut.

}
{} {}

\inputaufgabe
{6}
{

Berikan suatu
\definitionsverweis {pemetaan injektif}{}{}
\definitionsverweis {kontinu}{}{}
\maabbdisp {\varphi} {\R} {S^2
} {,}
yang
\zusatzklammer {sebagai pemetaan ke $\R^3$} {} {}
\definitionsverweis {dapat direktifikasi}{}{}
memiliki
\definitionsverweis {panjang}{}{}
tak hingga, dan memenuhi
\mathkor {} {\operatorname{lim}_{   t  \rightarrow  \infty  } \, \varphi (t) = N} {dan} {\operatorname{lim}_{   t  \rightarrow  - \infty  } \, \varphi (t) = S} {}.

}
{} {}

\inputaufgabe
{4}
{

Tunjukkan bahwa sumbu-sumbu koordinat bukan
\definitionsverweis {manifold topologis}{}{}.

}
{} {}
